% --- Back cover ---
% Large PERSONIX wordmark, no logo image. Cream paper inherited globally.
% Absolute (tikz overlay) positioning so the footer sticks to the bottom.
\clearpage
\thispagestyle{empty}
\begin{tikzpicture}[remember picture, overlay]
  % --- Wordmark ---
  \node[anchor=north, font=\sffamily\bfseries\fontsize{48}{52}\selectfont]
    at ([yshift=-26mm]current page.north) {PERSONIX};
  \node[anchor=north, font=\rmfamily\itshape\large]
    at ([yshift=-44mm]current page.north) {שינוי ללא פשרות};
  \node[anchor=north, font=\rmfamily\small,
        text width=\paperwidth-50mm, align=center]
    at ([yshift=-54mm]current page.north)
    {רשת מוניטין מבוזרת שאי אפשר לצנזר ואי אפשר לשחד};
  % --- Body ---
  \node[anchor=north, text width=\paperwidth-50mm, align=justify, font=\small]
    at ([yshift=-72mm]current page.north)
    {%
      המדינה עבדה כשהתקשורת הייתה אטית, ההחלטות היו מקומיות, והסמכות
      חיה באותה עיר שבה שלטה. רשתות ימינו הופכות את כל שלוש
      ההנחות --- והמוסדות הבנויים עליהן כבר אינם מתאימים
      לעולם שבו הם שולטים. הספר הזה מתאר כלי שכן מתאים:
      רשת מוניטין מבוזרת שבה הזהות שלך לשמור,
      האמת ניתנת לביקורת פומבית, והשוחד הוא התאבדות מבחינת המוניטין.

      \smallskip
      לא מניפסט. לא תחזית. שרטוט עבודה לכך כיצד אפשר לבנות
      את יורשת המדינה --- מרצון, הצהרה אחת
      בכל פעם.
    };
  % --- Footer (anchored to the bottom of the page) ---
  \node[anchor=south, font=\sffamily\footnotesize,
        text width=\paperwidth-50mm, align=center]
    at ([yshift=12mm]current page.south)
    {%
      Pavel Kudrna\quad\textbullet\quad פראג, 2026\par
      \href{https://personix.org}{personix.org}\par
      מופץ תחת רישיון Creative Commons Attribution-ShareAlike 4.0 International
      (CC BY-SA 4.0).
    };
\end{tikzpicture}%
\mbox{}%
\clearpage
